\chapter{Hơn thua với người}
\markboth{Hơn thua với người}{Competing with Others}

\section{Anh em trong nhà lớn lên bằng so sánh.}
\begin{truyen}{Anh em trong nhà lớn lên bằng so sánh.}{Comparison can wound love.}
\chuhoa{T}{ừ nhỏ, hai anh em đã bị đặt cạnh nhau: ai giỏi hơn, ai ngoan hơn, ai làm ra tiền nhiều hơn. Càng lớn họ càng xa. Không phải vì ghét nhau, mà vì mối quan hệ đã bị biến thành một cuộc thi quá lâu.}
\textit{(From childhood, two siblings were constantly compared: who was smarter, better, or earned more. The older they grew, the farther apart they became. Not because they hated each other, but because the relationship had been turned into a competition for too long.)}
\end{truyen}
\begin{baihoc}
Có những quan hệ đáng ra để nâng nhau, nhưng bị làm hỏng vì hơn thua kéo dài.
\textit{(Some relationships are meant to support each other, but are damaged by prolonged competition.)}
\end{baihoc}
\begin{ghinhoanh}\textbf{Short English to remember:}

Comparison can wound love.
\end{ghinhoanh}
\begin{tuhocnhanh}
1. Em có đang bị đặt trong một cuộc so sánh nào không? \textit{(Are you trapped in any comparison right now?)}

2. Em muốn bước ra khỏi nó bằng cách nào? \textit{(How do you want to step out of it?)}
\end{tuhocnhanh}
\ngancach

\section{Bạn bè gặp nhau mà chỉ nói chuyện ai hơn ai.}
\begin{truyen}{Bạn bè gặp nhau mà chỉ nói chuyện ai hơn ai.}{A friendship can become a scoreboard.}
\chuhoa{L}{âu ngày, nhóm bạn cũ không còn hỏi nhau có khỏe không, chỉ hỏi lương bao nhiêu, nhà mua chưa, con học trường gì. Gặp nhau nhiều mà ai cũng mệt. Khi mọi cuộc gặp thành bảng điểm, tình bạn bắt đầu mòn.}
\textit{(Over time, old friends stopped asking how each other truly felt and only asked about salary, houses, and schools. They met often yet everyone left tired. When every reunion becomes a scoreboard, friendship starts to wear out.)}
\end{truyen}
\begin{baihoc}
Tình bạn không sống được lâu nếu mỗi lần gặp nhau chỉ để đo hơn kém.
\textit{(Friendship cannot last long if every meeting is used to measure who is ahead.)}
\end{baihoc}
\begin{ghinhoanh}\textbf{Short English to remember:}

A friendship can become a scoreboard.
\end{ghinhoanh}
\begin{tuhocnhanh}
1. Trong các mối quan hệ, em có đang bị đo bằng thành tích không? \textit{(In your relationships, are you being measured mainly by achievements?)}

2. Em muốn giữ lại kiểu trò chuyện nào hơn? \textit{(What kind of conversation do you want to preserve instead?)}
\end{tuhocnhanh}
\ngancach

\section{Một người hơn đồng nghiệp một chút nhưng kém lòng rộng lượng rất nhiều.}
\begin{truyen}{Một người hơn đồng nghiệp một chút nhưng kém lòng rộng lượng rất nhiều.}{Being ahead is not the same as being better.}
\chuhoa{A}{nh làm giỏi hơn thật, nhưng luôn khó chịu khi thấy người khác tiến bộ. Anh hơn họ ở vài kỹ năng, nhưng thua họ ở chỗ biết chúc mừng và biết nâng nhau lên.}
\textit{(He truly performed better, yet always felt irritated when others improved. He was ahead in some skills, but behind in generosity and the ability to celebrate others.)}
\end{truyen}
\begin{baihoc}
Hơn người ở thành tích mà thua người ở khí độ thì cái hơn đó rất ngắn.
\textit{(If you are ahead in achievement but behind in character, the advantage is short-lived.)}
\end{baihoc}
\begin{ghinhoanh}\textbf{Short English to remember:}

Ahead is not always better.
\end{ghinhoanh}
\begin{tuhocnhanh}
1. Em có vui được thật lòng khi người khác khá lên không? \textit{(Can you sincerely feel glad when others improve?)}

2. Em đang muốn hơn ai ở điều gì? \textit{(Whom are you trying to be ahead of, and in what?)}
\end{tuhocnhanh}
\ngancach

\section{Một người cứ phải chứng minh mình đúng trước cha mẹ.}
\begin{truyen}{Một người cứ phải chứng minh mình đúng trước cha mẹ.}{Some victories do not heal family.}
\chuhoa{C}{ậu trưởng thành rồi, giỏi rồi, nên mỗi lần tranh luận với cha mẹ đều muốn chốt phần thắng. Nhưng sau nhiều lần như thế, nhà vẫn nặng. Cậu bắt đầu hiểu người thân không phải lúc nào cũng cần mình thắng lý, họ cần mình còn giữ cách nói có tình.}
\textit{(He grew up and became capable, so every discussion with his parents became a chance to prove he was right. After many such moments, the house still felt heavy. He began to understand that family does not always need logical victory, but speech with tenderness.)}
\end{truyen}
\begin{baihoc}
Trong gia đình, thắng bằng lý nhiều khi vẫn thua bằng không khí.
\textit{(In a family, you can win in logic and still lose in atmosphere.)}
\end{baihoc}
\begin{ghinhoanh}\textbf{Short English to remember:}

Some victories do not heal family.
\end{ghinhoanh}
\begin{tuhocnhanh}
1. Em có đang cố chứng minh nhiều hơn đang cố hiểu trong gia đình không? \textit{(Are you trying to prove more than understand at home?)}

2. Em có thể đổi cách nói nào để giữ không khí tốt hơn? \textit{(What way of speaking can you change to keep a better atmosphere?)}
\end{tuhocnhanh}
\ngancach

\section{Một cô gái không muốn thua người cũ nên sống rất mệt.}
\begin{truyen}{Một cô gái không muốn thua người cũ nên sống rất mệt.}{Living to beat someone is another kind of loss.}
\chuhoa{S}{au chia tay, cô không còn sống cho mình nữa mà sống để chứng minh cho ai đó thấy mình ổn, mình đẹp, mình thành công hơn. Càng cố hơn thua, cô càng xa khỏi đời sống thật của mình.}
\textit{(After a breakup, she no longer lived for herself but to prove to someone else that she was fine, attractive, and successful. The more she competed, the further she drifted from her real life.)}
\end{truyen}
\begin{baihoc}
Khi mục tiêu sống chỉ còn là hơn một người khác, ta đã giao tay lái đời mình cho họ.
\textit{(When life becomes mainly about being ahead of someone else, we hand them the steering wheel of our life.)}
\end{baihoc}
\begin{ghinhoanh}\textbf{Short English to remember:}

Living to beat someone is another loss.
\end{ghinhoanh}
\begin{tuhocnhanh}
1. Em có đang để ai điều khiển lựa chọn của mình bằng cảm giác hơn thua không? \textit{(Is anyone controlling your choices through rivalry?)}

2. Nếu sống lại cho mình, em sẽ bớt làm điều gì? \textit{(If you returned to living for yourself, what would you stop doing?)}
\end{tuhocnhanh}
\ngancach